%!TEX root = ../thesis.tex
\chapter{多机器人系统架构设计}
\label{chap:main}
相较于传统单机器人搜索模式，多机器人协同通过任务分配、资源共享与动态协调，能够有效提升搜索效率、扩大覆盖范围并增强系统容错能力。然而，多机器人协同搜索系统的实现面临诸多问题：机器人间的实时通信、动态环境下的任务分配优化、机器人碰撞规避以及系统可扩展性设计等问题，均需通过合理的系统架构予以支撑。本章基于机器人操作系统ros2-galactic的分布式框架，提出一种多机器人协同搜索系统架构，包括机器人部署模块，地图融合模块，目标检测模块，目标位置解算模块，路径规划模块。该架构旨在实现机器人资源的高效调度、环境信息的融合共享以及优化协同搜索行为，为后续章节中目标检测算法、路径规划策略及仿真实验验证提供基础框架。

\section{系统开发与运行平台}
\label{sec:font}
\subsection{机器人操作系统ROS}
ROS（Robot Operating System）是当下最受欢迎与
广泛使用的开源机器人操作系统。它源于斯坦福大学人工智能实验室 STAIR 项目创建
的软件系统原型，2007 年 Willow Garage 公司与之合作，在积累了大量研究人员的科学
研究成果后逐渐形成了 ROS 核心概念和基础软件架构。ROS 的分布式架构、模块化设计及丰富的工具链，为多机器人协同系统提供了高效开发基础。ROS 具有以下几个特性：

\begin{itemize}
\item 开源性与社区生态

ROS 遵循 BSD 开源协议，允许个人及企业免费使用与二次开发。其庞大的社区生态（如 ROS Wiki、GitHub 仓库）提供了海量开源功能包，涵盖传感器驱动、SLAM、路径规划等核心模块，极大降低了机器人系统的开发门槛。
\item 分布式架构与松耦合设计

ROS系统的通信架构允许建立多个节点，通过话题和订阅等方式来管理机器人通信，十分适用于模块化开发，有利于调整增加机器人的功能，使系统具有很好的扩展性和鲁棒性。

\item 多语言与跨平台支持

ROS 支持 C++、Python 等主流编程语言，其中 Python 语法简洁，软件包支持丰富，具有开发效率上的优势。ROS2 Galactic 版本全面适配 Python 3.8+，并强化了对嵌入式平台（如树莓派）、实时操作系统（RTOS）及 Windows 系统的兼容性。
\item 第三方工具与算法集成

ROS 可无缝集成 OpenCV、PCL（点云库）、TensorFlow 等第三方库，同时支持 Gazebo 仿真环境、RViz 可视化工具及 Qt 人机界面框架。ROS2 Galactic 进一步优化了与 DDS 中间件（如 Cyclone DDS）的兼容性，为多机器人协同提供低延迟通信保障
\item 选用ROS2 Galactic

尽管 ROS1 在单机器人系统中表现优异，但其通信实时性不足、网络依赖性强等缺陷限制了多机器人协同场景的应用。ROS2 Galactic 作为长期支持（LTS）版本，针对上述问题进行了全面革新：
  \begin{itemize}
    \item 去中心化通信：基于 DDS 协议实现多机器人动态组网，避免单点故障，支持定制 QoS 策略（如低延迟、高可靠性传输）;

    \item 实时性与安全性：通过生命周期节点管理、线程安全接口及加密通信机制，提升系统稳定性；

    \item Python 开发优势：ROS2 Galactic 提供成熟的 rclpy 库，结合 Python 的轻量级代码与丰富算法库（如 NumPy、PyTorch），可快速实现任务分配、动态路径规划等协同逻辑，显著提升开发效率。
  \end{itemize}
综上，ROS2 Galactic 凭借其分布式架构、强实时通信及对 Python 生态的深度支持，成为构建高可靠、可扩展多机器人协同搜索系统的理想技术基座。
\end{itemize}

\subsection{Gazebo 机器人仿真平台}
Gazebo 是一款机器人仿真软件平台，具有高性能的 ODE 物理引擎，提供高质量的
三维图形并支持多种编程语言。Gazebo 支持多种机器人传感器模型（如激光传感器、
IMU、Kinect 等），提供多种机器人模型（如 PR2、Turtlebot、机械臂等）并支持定制
私人机器人模型，并且能够对真实环境中物理实体进行建模，为开发人员编写机器人应
用程序以及算法快速测试提供了良好的仿真平台。
\begin{figure}[htbp] % use float package if you want it here
  \centering
  \includegraphics[width=\textwidth]{TheGazebo.png}
  \caption{Gazebo 机器人仿真平台}
  \label{fig:ipe}
\end{figure}

\subsection{移动机器人Turtlebot3}
移动机器人是多机器人系统中最重要的执行本体，是一个集环境感知、任务规划、
动作执行等功能于一体的硬件系统。TurtleBot3 是第三代基于ROS的开源移动机器人平台，由 ROBOTIS 公司设计并开源，旨在为教育、科研、产品原型开发提供低成本、高灵活性的机器人开发工具。TurtleBot3 的尺寸仅为前代产品的 1/4（如 Burger 型号尺寸为 138×178×192mm），重量仅 1kg，便于携带。其价格远低于同类 SLAM 机器人，适合教育场景的规模化应用。并且作为 ROS 官方支持的机器人平台，TurtleBot3 原生兼容 ROS 的导航栈（如 Gmapping、Cartographer）和工具链（如 Gazebo 仿真、RViz 可视化），使用效率更高。

TurtleBot3 提供 Burger 和 Waffle Pi 两种主流型号，主要区别如下：
\begin{table}[ht]
  \centering
    \begin{tabular}{lcc}
      \toprule
      机器人型号 & TurtleBot3 Burger & TurtleBot3 Waffle \\
      \midrule
      最大平移速度   & 0.22 m/s                   & 0.26 m/s                   \\
      最大角速度     & 2.84 rad/s                 & 1.82 rad/s                 \\
      最大负载       & 15 kg                      & 30 kg                      \\
      尺寸 (长$\times$宽$\times$高) & 138 mm $\times$ 178 mm $\times$ 192 mm & 281 mm $\times$ 306 mm $\times$ 141 mm \\
      自重           & $\sim$1 kg                 & $\sim$1.8 kg               \\
      预计工作时间   & 2.5 小时                   & 2 小时                     \\
      控制单元       & 树莓派 3                  & Intel Joule 570x           \\
      传感器配置     & 360°激光雷达 (LDS-01)；无摄像头 & 360°激光雷达 (LDS-01)；内置摄像头\\
      \bottomrule
    \end{tabular}
  \caption{TurtleBot3 Burger 与 Waffle 型号对比}
  \label{tab:turtlebot3_comparison}
\end{table}

为了提高环境感知能力，同时允许机器人通过光学摄像头识别搜索目标，本项目选用Turtlebot3 Waffle机器人。

\section{总体架构与流程}

多机器人目标搜索系统利用 ROS2 将各个资源与模块相连接，使得各个节点之间具有松耦合和模块化的特性，保证了系统的可扩展性、稳定性和灵活性，系统工作的总体流程示意图如图2-2所示。

每个ROS2包分别负责不同的功能，包括机器人控制、仿真环境搭建、目标检测等，系统主要由以下ROS2包构成,根据功能构成多个模块：

\begin{table}[ht]
  \centering
    \begin{tabular}{lc}
      \toprule
      ROS包名 & 主要功能\\
      \midrule
      merge\_map   & 地图融合                               \\
      multi\_robot\_exploration     & 在仿真环境中部署机器人         \\
      path\_follow       & 根据规划路径控制turtlebot机器人的运动                                     \\
      target\_detect & 目标搜索与坐标解算  \\
      turtlebot3\_gazebo       &   仿真环境搭建与配置             \\
      yolov5\_ros2   &   yolo模型目标识别          \\
      \bottomrule
    \end{tabular}
  \caption{系统主要ROS包及其功能}
  \label{tab:turtlebot3_comparison}
\end{table}

\begin{figure}[ht] % use float package if you want it here
  \centering
  \includegraphics[width=\textwidth]{流程图1.png}
  \caption{系统基本架构图}
  \label{fig:ipe}
\end{figure}

\subsection{地图融合模块}
全局地图融合是多机器人协作探索环境和地图构建系统中的关键部分之一。全局地
图融合模块首先加载多机器人初始相对位姿参数，实时检测多机器人局部地图信息并对
地图数据接收。根据当前系统中的相关参数以及多机器人实时构建的局部地图确定全局
地图的尺寸与位姿，最后根据多张局部地图信息进行全局地图融合。

\subsection{路径规划模块}
路径规划模块将根据机器人SLAM得到的地图进行障碍地图的生成，并基于此进行路径规划以完成对未知环境的探索。使用了A*等算法进行动态路径规划，可以同时为多个机器人规划路径，为了适用与多机器人搜索的要求，路径规划时会考虑多机器人之间的避碰，并且提高搜索效率。规划模块会将路径话题发布到包path\_follow进行解算后最终控制机器人的运动。

\subsection{目标检测模块}
目标检测模块使用了机器人上的摄像头模块，利用yolov5s人工智能视觉模型对目标进行识别，并且将识别的数据与激光雷达的距离数据进行融合，最后结合机器人位姿坐标进行目标坐标的标记与解算。为了提高效率与精确度，传感器数据融合过程使用了时间同步策略，并且对解算得到的原始坐标进行后期数据处理，使用聚类分析等数据处理技术对目标坐标进行推理优化。

